% src: PREREGISTRATION.md; PREREGISTRATION_ROUND2.md; DEVIATIONS.md DEV-8; revision/paper/main.tex
\section{Discussion and Limitations}
\label{sec:discussion}
\subsection{What a transferred operating point establishes}
The round-1 audit distinguishes a detector that orders target scores poorly from one whose source threshold sits
poorly on otherwise informative scores. Both can produce a high deployed error, but they suggest different responses.
Changing a threshold cannot repair poor discrimination. Conversely, selecting a detector solely by target-oracle EER
does not address how a source threshold will behave on genuine target speech. The false-rejection and
false-acceptance decomposition is therefore needed alongside HTER and its transfer gap.

% src: PREREGISTRATION.md; PREREGISTRATION_ROUND2.md; DEVIATIONS.md DEV-8; revision/paper/main.tex
The reporting card is intended to make these distinctions inspectable. It requires the threshold source and
estimator, the application policy, the target error components and an oracle reference computed only for analysis.
The target oracle is not an available deployment policy unless labelled target data can actually be collected.
Similarly, the genuine-speech-only threshold reset in round~1 measures a tradeoff: restoring a rejection budget
does not establish a useful spoof detector. A workflow that uses the scores for review must assess review load and
automatic errors separately; Section~\ref{sec:r2review} shows that a review budget designed on development data is not
kept after transfer.

\subsection{What carries over from round 1}
% src: results/round2/analysis/r2_registered.txt; results/round2/analysis/r2_explore.txt; gen/r2_operating.tex; gen/r2s_last.tex
Round~2 tested the round-1 observations on new scores, and four of them carry over. First, the dominant error on a
shifted target follows the operating policy (R2-H1, two of three detectors): a strict false-rejection budget puts the
error on genuine ITW speech, and a lenient false-acceptance budget puts it on spoofs. The exception, the ASVspoof~5
baseline, rejects almost all genuine ITW speech even at its most lenient registered threshold; there the target shift
overrides the policy. Second, the transfer gap on LA19 eval belongs to the AASIST-L recipe and its checkpoint selection,
not only to the released file (R2-H2), although its size depends on the seed and on the selection rule. Third, better
discrimination did not bring better transfer: both new LA19-trained detectors discriminate ITW far better than the
round-1 primary models and are transfer-limited, as the round-1 reference was. Fourth, band-limiting and low-bitrate
channels raise false rejection: R2-H4 holds for both low-bitrate speech codecs and the phone impulse response, and the
narrowband conditions of ASVspoof~5 show the same pattern, whereas the effect of additive noise depends on the model.

% src: results/round2/analysis/r2_registered.txt; gen/r2_seeds.tex; gen/r2s_last.tex
Two observations do not carry over in the form that round~1 suggested. Removing silence from training lowers ITW false
rejection at the development threshold, but about two thirds of that reduction comes from checkpoint selection, FAR
rises, and ITW discrimination does not improve. The silence cue thus helps explain where this recipe's ITW operating
point falls, not why the recipe fails to discriminate ITW. And a released checkpoint's ITW operating point is not the
recipe's: the released AASIST-L rejects more genuine ITW speech than any of its retrained counterparts.

% src: results/round2/analysis/r2_registered.txt; results/round2/analysis/r2_explore.txt
For reporting, round~2 strengthens three items of the card (Table~\ref{tab:card}). The threshold policy and its source
belong with every rate, because they decide which error dominates. The dev-equivalent interval belongs with every
\(\tauE\) set on nearly separable development data: across that interval, XLSR-Mamba's ITW FRR spans 31.4\,pp. And a
review policy has to be assessed by its review share and its automatic errors on the target: a budget set on development
data is not kept after transfer, and re-anchoring it without labels keeps the budget only by automating errors.

\subsection{Scope of inference}
The released-checkpoint audit inherits each author's training, augmentation and model-selection choices. Its roster
was selected for a broader deployment project rather than sampled from the population of speech detectors. Related
models share a lineage, and the self-supervised reference enters no primary-model hypothesis count. The registered
extension broadens this evidence but still retrains only AASIST-L. An observation across its seeds supports a
statement about that recipe under the recorded training and selection rules; it does not identify an architecture
effect across independently trained detector families.

% src: PREREGISTRATION.md; PREREGISTRATION_ROUND2.md; DEVIATIONS.md DEV-8; revision/paper/main.tex
The ASVspoof~5 detector differs from LA19-trained detectors in more than silence handling, including its training
corpus. Its training audio was not subject to the development/evaluation post-processing
\citep{wang25_asvspoof5_db}. It therefore cannot isolate a causal effect of removing silence during training. The
trimmed and untrimmed AASIST-L arms provide the registered intervention. Even there, the intervention changes both
training and threshold-source preprocessing, and the recipe's weight-averaging eligibility did differ between arms: no
trimmed run reached the averaging threshold. The selection-free last-epoch comparison separates most of that difference,
but it is exploratory, and five runs per arm limit any permutation test to \(p\geq1/252\). Test-time trimming alone, as
in round~1, can expose reliance on a cue but cannot establish why a detector learned it.

% src: PREREGISTRATION.md; PREREGISTRATION_ROUND2.md; DEVIATIONS.md DEV-8; revision/paper/main.tex
\subsection{Data, dependence and ecological validity}
ITW is a community-exposed benchmark, not a prospective sample of previously unseen deployment traffic
\citep{muller22_in_the_wild}. Its metadata identify speakers but not source recordings. Speaker-cluster inference
accounts for within-speaker dependence under the recorded labels; recording-specific effects and dependence that
crosses speaker labels cannot be separated from those labels. The disagreement between clip-level and
speaker-level H2 inference is part of the audit's evidence, rather than a reason to substitute one registered
result for the other. The exploratory H1 and H4 comparisons likewise rest on few distinct models, so multiple
model-policy points do not create independent model replications.

% src: PREREGISTRATION.md; PREREGISTRATION_ROUND2.md; DEVIATIONS.md DEV-8; revision/paper/main.tex
The DF21 analysis uses the documented integrity-verified evaluation surface. Excluding audio that cannot reproduce
its own encoded checksum changes the evaluated population, and the clean-surface sensitivity analysis does not
recover the omitted trials. Our DF21 measurements must therefore be identified by their surface and should not be
equated with published measurements on a different subset or decode. The checksum verification and decoder
amendment are reproducibility requirements, not an estimate that the excluded trials are harmless.

% src: PREREGISTRATION.md; PREREGISTRATION_ROUND2.md; DEVIATIONS.md DEV-8; revision/paper/main.tex
MLAAD paired with M-AILABS adds languages and synthesis systems but also a class-associated sample-rate difference.
Resampling traces may contribute to separation, so a pooled or language-specific error rate cannot be interpreted
as an isolated language effect. The original round-1 MLAAD scores were already known when the extension was
registered; their reporting remains descriptive. Bona fide speaker and spoof-system clustering measures the
available dependence structure, not independence of every utterance.

% src: PREREGISTRATION.md; PREREGISTRATION_ROUND2.md; DEVIATIONS.md DEV-8; revision/paper/main.tex
The stress experiments preserve clip identity and isolate specified transformations. A codec reconstruction of
genuine speech retains its genuine label under this channel-use interpretation; codec-based generation can have
another role in synthesis benchmarks \citep{wu24_codecfake,wang25_asvspoof5_db,xiao26_resynth}. The experiments cover
single-pass digital transformations and impulse-response convolution. They do not reproduce an entire live
transmission chain, a physical re-recording campaign or an adaptive adversary. Desktop scoring does not establish
on-device numerical parity, latency or energy performance.

% src: PREREGISTRATION.md; PREREGISTRATION_ROUND2.md; DEVIATIONS.md DEV-8; revision/paper/main.tex
\subsection{Registration and reporting boundaries}
The tagged registration and deviation log provide a version-controlled chronology rather than an independently
deposited registration. Some original amendments followed pilot or reproduction observations, and the speaker-level
sensitivity analyses were requested after the round-1 results were known. The manuscript reports that chronology,
preserves the original clip-level outcomes and labels the added analyses as exploratory. Round~2 was designed in
response to those observations; its registration protects the later tests from adaptation to its own scores,
without making the original observations unknown at design time.

Our costs are explicit evaluation scenarios, not measured operational priors or monetary losses. A review policy
does not model reviewer accuracy, staffing, delay or consent requirements for collecting genuine channel speech.
The audit characterizes detector limitations and the evidence needed to assess them; it does not certify an
authenticity decision for a particular speaker or forensic case.
